\documentclass[12pt]{article}
\usepackage[a4paper, margin=2cm]{geometry}

\usepackage{fontspec}
\usepackage{polyglossia}
\setdefaultlanguage{russian}
\setmainlanguage{russian}
\setmainfont{Times New Roman}
\begin{document}

\begin{center}
\textbf{\LARGE ТЕХНИЧЕСКИЙ ОТЧЕТ}
\end{center}

\vspace{1cm}

% === ТЕКСТ ДОКУМЕНТА ВСТАВЛЯТЬ СЮДА ===
% Здесь должно быть описание проверяемого оборудования,
% проведенные измерения, анализ состояния и выводы.
% Допускается ссылка на нормативные документы и регламенты.

\noindent
\textbf{Тема:} Технический отчет по изменению котировок (динамике) акций наиболее популярных нефтедобывающих компаний за последний месяц по состоянию на 28.03.2026.

\bigskip

\noindent
\textbf{1. Цель и область применения.} 
Цель настоящего отчета — подготовить сводную техническую оценку динамики котировок акций крупнейших (наиболее популярных) нефтедобывающих компаний за период последнего месяца, приуроченного к дате \textbf{28.03.2026}. Отчет предназначен для аналитического сопровождения процесса принятия решений и включает:
\begin{itemize}
    \item определение перечня эмитентов (нефтедобывающих компаний), подлежащих оценке;
    \item описание методики расчета изменений котировок за период;
    \item представление качественных выводов о характере изменения котировок (рост/снижение/волатильность);
    \item формирование технических ограничений и оговорок по качеству данных.
\end{itemize}

\bigskip

\noindent
\textbf{2. Исходные данные и ограничения.}
В рамках формирования отчета использованы результаты автоматизированного запроса к корпоративной базе знаний и веб-поиска по следующим направлениям:
(а) котировки/динамика акций нефтедобывающих компаний за период последнего месяца по состоянию на 28.03.2026; 
(б) перечень крупнейших нефтедобывающих компаний и их тикеров; 
(в) корпоративная аналитика/отчетность по рынку акций нефтегазового сектора за последний месяц.

\medskip
\noindent
\textbf{Примечание по данным:} в доступном контексте отсутствуют конкретные числовые значения котировок, тикеры и ссылки на внутренние аналитические материалы (контент для извлечения помечен как \emph{N/A}). В связи с этим в отчете применяется нейтральный заполнитель полей анализа, сохраняя формальную структуру и техническую логику документа. При предоставлении фактических данных (цены закрытия/диапазон котировок/объемы) разделы \textbf{4} и \textbf{5} подлежат заполнению расчетными показателями.

\bigskip

\noindent
\textbf{3. Перечень эмитентов и критерии отбора.}
\begin{itemize}
    \item \textbf{Критерий «популярные/крупнейшие нефтедобывающие компании»:} включает эмитентов нефтедобывающего профиля с высокой ликвидностью торгов и широким информационным присутствием (конкретный перечень и тикеры в данном документе не верифицированы из-за отсутствия данных в контексте).
    \item \textbf{Период оценки:} один месяц, предшествующий \textbf{28.03.2026} (конкретная стартовая дата периода обозначается как \textbf{T-1 месяц}).
\end{itemize}

\medskip
\noindent
\textbf{Таблица 1.} Перечень компаний (заполнить после получения тикеров и доступных котировок).
\begin{itemize}
    \item Компания 1 (тикер: \rule{4cm}{0.4pt}) — нефтедобывающий профиль
    \item Компания 2 (тикер: \rule{4cm}{0.4pt}) — нефтедобывающий профиль
    \item Компания 3 (тикер: \rule{4cm}{0.4pt}) — нефтедобывающий профиль
    \item Компания 4 (тикер: \rule{4cm}{0.4pt}) — нефтедобывающий профиль
    \item Компания 5 (тикер: \rule{4cm}{0.4pt}) — нефтедобывающий профиль
\end{itemize}

\bigskip

\noindent
\textbf{4. Методика расчета динамики котировок.}
Для каждого эмитента рассматриваются котировки на опорные даты:
\begin{itemize}
    \item $P_{start}$ — цена закрытия (или согласованная базовая котировка) на дату \textbf{T-1 месяц};
    \item $P_{end}$ — цена закрытия (или базовая котировка) по состоянию на \textbf{28.03.2026}.
\end{itemize}

\medskip
\noindent
\textbf{Ключевые показатели (для последующего заполнения):}
\begin{itemize}
    \item \textbf{Абсолютное изменение:} $\Delta P = P_{end} - P_{start}$;
    \item \textbf{Процентное изменение:} $\Delta \% = \dfrac{P_{end} - P_{start}}{P_{start}} \cdot 100\%$;
    \item \textbf{Оценка волатильности за период:} по согласованным статистикам (например, стандартное отклонение дневных доходностей или диапазон $High/Low$ за период).
\end{itemize}

\medskip
\noindent
\textbf{Допущения:} 
используется единая методика расчета для всех эмитентов (одинаковый тип цены: закрытие/средняя/последняя доступная; одинаковая торговая площадка/класс акций). При отсутствии единообразия исходных параметров значения подлежат корректировке либо указываются отдельными подпунктами.

\bigskip

\noindent
\textbf{5. Результаты технической оценки (заполнить после получения данных).}
\medskip
\noindent
\textbf{Таблица 2.} Динамика котировок за период (заполнить фактическими значениями).
\begin{itemize}
    \item Компания 1 (тикер: \rule{4cm}{0.4pt}): $P_{start}=\rule{3cm}{0.4pt}$, $P_{end}=\rule{3cm}{0.4pt}$, $\Delta P=\rule{3cm}{0.4pt}$, $\Delta \%=\rule{2.5cm}{0.4pt}$, волатильность: \rule{3.5cm}{0.4pt}
    \item Компания 2 (тикер: \rule{4cm}{0.4pt}): $P_{start}=\rule{3cm}{0.4pt}$, $P_{end}=\rule{3cm}{0.4pt}$, $\Delta P=\rule{3cm}{0.4pt}$, $\Delta \%=\rule{2.5cm}{0.4pt}$, волатильность: \rule{3.5cm}{0.4pt}
    \item Компания 3 (тикер: \rule{4cm}{0.4pt}): $P_{start}=\rule{3cm}{0.4pt}$, $P_{end}=\rule{3cm}{0.4pt}$, $\Delta P=\rule{3cm}{0.4pt}$, $\Delta \%=\rule{2.5cm}{0.4pt}$, волатильность: \rule{3.5cm}{0.4pt}
    \item Компания 4 (тикер: \rule{4cm}{0.4pt}): $P_{start}=\rule{3cm}{0.4pt}$, $P_{end}=\rule{3cm}{0.4pt}$, $\Delta P=\rule{3cm}{0.4pt}$, $\Delta \%=\rule{2.5cm}{0.4pt}$, волатильность: \rule{3.5cm}{0.4pt}
    \item Компания 5 (тикер: \rule{4cm}{0.4pt}): $P_{start}=\rule{3cm}{0.4pt}$, $P_{end}=\rule{3cm}{0.4pt}$, $\Delta P=\rule{3cm}{0.4pt}$, $\Delta \%=\rule{2.5cm}{0.4pt}$, волатильность: \rule{3.5cm}{0.4pt}
\end{itemize}

\medskip
\noindent
\textbf{Качественная характеристика итогов (пока без числовых данных):}
\begin{itemize}
    \item по группе эмитентов в целом ожидается выявление доминирующего тренда (рост/снижение) по знаку $\Delta \%$;
    \item уровень волатильности подлежит интерпретации как индикатор повышенной неопределенности на рынке в течение месяца;
    \item при наличии разнонаправленной динамики выделяется подгруппа компаний с наибольшими абсолютными отклонениями $\Delta P$ и $\Delta \%$.
\end{itemize}

\bigskip

\noindent
\textbf{6. Аналитические выводы.}
На дату подготовки отчета количественные результаты динамики котировок не могут быть представлены в числовом виде из-за отсутствия в доступном контексте необходимых исходных данных (котировки, тикеры, аналитические выгрузки). После предоставления данных отчет подлежит актуализации с внесением расчетных значений $\Delta P$, $\Delta \%$ и показателей волатильности для каждого эмитента.

\medskip
\noindent
\textbf{Ожидаемые направления использования результатов:}
\begin{itemize}
    \item формирование краткосрочного рыночного профиля нефтедобывающего сектора на интервале \textbf{T-1 месяц} — \textbf{28.03.2026};
    \item поддержка аналитических сценариев (рост/стагнация/снижение) и оценка чувствительности к рыночным факторам;
    \item обоснование дальнейшего мониторинга на регулярной основе.
\end{itemize}

\bigskip

\noindent
\textbf{7. Технические оговорки.}
\begin{itemize}
    \item Отчет носит технический характер и отражает методику оценки; фактическая численная интерпретация требует корректного ввода исходных котировок.
    \item При использовании данных из внешних источников необходимо обеспечить их сопоставимость (площадка, инструмент, тип цены, режим торгов).
    \item При необходимости могут быть добавлены ссылки на внутренние регламенты мониторинга финансовых инструментов и на применимые нормативные акты/стандарты раскрытия информации (в настоящем документе ссылки не приводятся ввиду отсутствия конкретных регламентов в контексте).
\end{itemize}

\vfill

\noindent
Дата: \rule{5cm}{0.4pt} \hfill Подпись: \rule{5cm}{0.4pt}

\end{document}